

\begin{table*}[htbp!]
\small
    \centering
    \begin{tabular}{lcccccccc} 
    \toprule
     & \multicolumn{1}{c}{} & \multicolumn{4}{c}{\textbf{Domain Specific Simulators}} & \multicolumn{3}{c}{\textbf{Analytical Models}} \\ 
    \cmidrule(l{2pt}r{2pt}){3-6}\cmidrule(l{2pt}r{2pt}){7-9}
    & \textbf{\sysname} & Proteus & vTrain & DistSim & Daydream & Calculon & AMPed & DistIR \\
    \noalign{\hrule height 0.5pt}
    \rowcolor{gray!15}
         \multicolumn{9}{c}{\textit{System Properties}}\\\hline Deployment-Free Prediction  & \greencheck & \greencheck & \redcross & \redcross & \redcross & \greencheck & \greencheck & \greencheck \\
Transparent -- \textit{No Code Modifications} & \greencheck & \redcross & \redcross & \redcross & \redcross & \redcross & \redcross & \redcross \\
Workload Agnostic & \greencheck & \greencheck & \redcross & \greencheck & \greencheck & \redcross & \redcross & \greencheck \\
\noalign{\hrule height 0.5pt}
    \rowcolor{gray!15}
         \multicolumn{9}{c}{\textit{Modeling Domain}}\\\hline
    Data Parallel & \greencheck & \greencheck & \greencheck & \greencheck & \greencheck & \greencheck & \greencheck & \greencheck \\
    Tensor Parallel & \greencheck & \greencheck & \greencheck & \greencheck & \redcross & \greencheck & \greencheck & \greencheck \\
     Pipeline Parallel & \greencheck & \greencheck & \greencheck & \greencheck & \redcross & \greencheck & \greencheck & \greencheck \\
    Sequence Parallel & \greencheck & \redcross & \redcross & \redcross & \redcross & \greencheck & \redcross & \redcross \\
    Pipeline Interleaving & \greencheck & \greencheck & \redcross & \redcross & \redcross & \greencheck & \redcross & \greencheck \\
    Distributed Optimizer & \greencheck & \greencheck & \redcross & \redcross & \redcross & \greencheck & \redcross & \greencheck \\
    Activation Recomputation & \greencheck & \greencheck & \redcross & \redcross & \redcross & \greencheck & \redcross & \redcross \\
    Gradient Accumulation & \greencheck & \redcross & \redcross & \redcross & \redcross & \greencheck & \redcross & \greencheck \\

    \bottomrule
    \end{tabular}
    \caption{
    Comparison of \sysname with existing performance modeling approaches. \sysname uniquely combines deployment-free prediction, transparency, and workload agnosticism, supporting a broad range of parallelism and optimization strategies across training configurations. Competing systems are limited in either coverage and often require code modifications.}
    \label{tab:features_systems}
\end{table*}



\section{Background}
\label{sec:background}

Deep learning training (DLT) workloads have grown to unprecedented scales, with state-of-the-art models now containing billions of parameters. For instance, training the Llama-3 405B required an estimated 864 exaflop-days of compute \cite{dubey2024llama}. Training such large-scale deep learning models requires parallelizing computation across hundreds or thousands of accelerator devices. This parallelization employs techniques such as Data Parallelism (DP), which replicates the model across devices; Tensor Parallelism (TP), which partitions individual layers; and Pipeline Parallelism (PP), which splits the model into stages. However, parallelization alone is insufficient; achieving high efficiency requires carefully balancing compute, memory, and communication bottlenecks.

\vheading{Balancing Resource Utilization} To address these bottlenecks, researchers have proposed various techniques that trade off different resources (Table \ref{tab:knob_impacts}). For instance, tensor parallelism can reduce memory pressure by partitioning layers across devices, but increases communication overhead due to frequent all-reduce operations between partitioned layers. Pipeline parallelism \cite{gpipe,pipedream} introduces pipeline bubbles that reduce compute efficiency, but enables parallelization at a comparatively low communication cost and memory pressure. The Zero Redundancy Optimizer (ZeRO) \cite{zero} shards model parameters, gradients, and optimizer states across workers, reducing memory pressure at the cost of increased communication. Activation checkpointing \cite{shah2020memory} performs additional recomputation to reduce memory usage, dropping activations after the forward pass.

Efficient hardware utilization requires careful composition of these techniques based on the model architecture, training parameters, and available resources. Each technique introduces additional tunable parameters that affect this balance. For example, the interleaved 1F1B pipelining schedule \cite{pipedream} reduces pipeline bubbles by assigning multiple micro-batches to each pipeline stage, but requires careful tuning of micro-batch counts to balance communication overlap. Similarly, ZeRO offers different sharding stages that provide varying tradeoffs between memory and communication.

\vheading{Composing Training Recipes} The vast configuration space generated by these techniques, combined with their interdependencies, makes tuning DLT workloads challenging. \Cref{fig:motivation:strongscalingconfigs} illustrates optimal configurations for training the same model with varying numbers of accelerators. With limited resources (16 devices), a combination of tensor and pipeline parallelism alleviates memory pressure. When scaling to 128 devices, the reduced per-device memory requirement allows leveraging data parallelism instead of tensor parallelism, avoiding additional communication overhead. Our experiments show that using the optimal configuration for 16 devices results in a 1.74\myx higher cost when applied to 128 devices, compared to the optimal configuration. This performance sensitivity makes it impractical to rely solely on heuristics or previous experience for configuration selection.

\textit{\textbf{Takeaway:} DLT workloads employ a rich set of parallelization and optimization techniques with unique resource tradeoffs. These techniques must be carefully composed to maximize hardware utilization.}


\begin{table}[htbp]
    \centering
    \begin{tabular}{@{\hspace{1px}}l@{\hspace{0px}}c@{\hspace{4px}}c@{\hspace{4px}}c@{\hspace{1px}}} 
    \toprule
    \textbf{Resource Load} & Compute & Memory & Network \\
    \midrule
    Data Parallel & \reddownarrow & \greendownarrow & \reduparrow \\
    Tensor Parallel & \reddownarrow & \greendownarrow & \reduparrow \\
    Pipeline Parallel & \reddownarrow & \greendownarrow & \reduparrow  \\
    Sequence Parallel & \reddownarrow & \greendownarrow & \reduparrow \\
    Pipeline Interleaving & \greenuparrow & \greendownarrow & \reduparrow \\
    Distributed Optimizer & \textbf{--} & \greendownarrow & \reduparrow \\
    Activation Recomputation & \reddownarrow & \greendownarrow & \textbf{--} \\
    Gradient Accumulation & \reddownarrow & \greendownarrow & \greendownarrow \\
    \bottomrule
    \end{tabular}
    \caption{Effect of configuration knobs on compute utilization, memory load, and network load in a fixed global batch size setting. This table highlights the trade-offs associated with each knob: while some configurations increase compute utilization, others may reduce memory or network load, illustrating the balancing act required to optimize  large-scale training jobs. \vspace{-1.5em}
    }
    \label{tab:knob_impacts}
\end{table}



Given these challenges, we pose the following question: \textit{Can we transparently, accurately, and efficiently predict the performance of arbitrary DLT configurations without access to target hardware?} Answering this question is crucial for enabling rapid exploration of the configuration space to identify resource-efficient training recipes.







































































